\section*{Hoofdstuk \ref{ch:b2:biogeochemical-cycles} --- Biogeochemische kringlopen}

\begin{solution}{exo:b2:biogeochemical-cycles:1}
Een \omterm{def:b2:biogeochemical-cycles:reservoir}{reservoir} is een voorraad van het element (een massa); een \omterm{def:b2:biogeochemical-cycles:reservoir}{flux} is een
overdrachtssnelheid tussen \omterm{def:b2:biogeochemical-cycles:reservoir}{reservoirs} (massa per jaar); de \omterm{def:b2:biogeochemical-cycles:reservoir}{verblijftijd} is
de massa van het \omterm{def:b2:biogeochemical-cycles:reservoir}{reservoir} gedeeld door zijn totale uitstroom, de gemiddelde
tijd die een atoom erin doorbrengt. Landplanten: $450/120 \approx 4$ jaar
(ruwweg de koolstof van een blad; \omterm{def:b2:plant-meristems:wood}{hout} blijft decennia, en het gemiddelde
verbergt de spreiding). Diepe oceaan: $37000/100 \approx 400$ jaar --- de tijd
voordat een uitgezakt molecuul het oppervlak weer ziet.
\end{solution}

\begin{solution}{exo:b2:biogeochemical-cycles:2}
$2.5\times 2.13 = \qty{5.3}{GtC}$ per jaar. $10/2.13 = \qty{4.7}{ppm}$
--- als alles in de lucht bleef; bij de \omterm{prop:b2:biogeochemical-cycles:carbon}{atmosferische fractie} van
\qty{45}{\%}, \qty{2.1}{ppm}.
\end{solution}

\begin{solution}{exo:b2:biogeochemical-cycles:3}
Openen: biologische binding door nitrogenase (cyanobacteriën,
rhizobia, \emph{Azotobacter}, \emph{Frankia}) en bliksem. Teruggeven:
\omterm{def:b2:microbial-metabolism:anaerobic}{denitrificatie} (facultatief anaerobe organismen zoals \emph{Pseudomonas}
en \emph{Paracoccus} in anoxische bodem en sediment) en anammox
(planctomyceten). Daartussen: ammonificatie van organische stikstof tot
ammonium door afbrekers (schimmels, bacteriën); nitrificatie van
ammonium tot nitriet (\emph{Nitrosomonas}, ammoniakoxiderende
archaea) en van nitriet tot nitraat (\emph{Nitrobacter}); en
assimilatie van ammonium of nitraat door planten en microben.
\end{solution}

\begin{solution}{exo:b2:biogeochemical-cycles:4}
Over geologische tijd kan stikstof altijd door binders worden gemaakt uit het
onuitputtelijke $\mathrm{N_2}$ van de lucht, en die binders worden begunstigd
zodra stikstof schaars is; fosfor heeft zo'n \omterm{def:b2:biogeochemical-cycles:reservoir}{reservoir} niet en
de aanvoer ervan ligt vast door de verwering, zodat het plafond voor de
productiviteit op lange termijn fosfor is. Binnen een seizoen is binding traag
en kostbaar, en de meeste planten kunnen het niet, zodat de beschikbare
stikstof hier en nu de groei beperkt --- en daarom is kunstmest
voornamelijk stikstof.
\end{solution}

\begin{solution}{exo:b2:biogeochemical-cycles:5}
$\tau = 1/0.2 = 5$ jaar; $M^{*} = F\tau = \qty{10}{t}$, dat is
$\qty{1}{g/m^3} = \qty{1000}{mg/m^3}$, sterk eutroof. Tijd tot
\qty{90}{\%}: $\tau\ln 10 = 11.5$ jaar. Halvering van de aanvoer halveert
de \omterm{def:b2:biogeochemical-cycles:reservoir}{stationaire toestand}, tot \qty{500}{mg/m^3}, bereikt op dezelfde
tijdschaal.
\end{solution}

\begin{solution}{exo:b2:biogeochemical-cycles:6}
Jaren 1960: \qty{0.9}{ppm/yr}, \qty{1.9}{GtC/yr}. Jaren 2010: \qty{2.3}{ppm/yr},
\qty{5.0}{GtC/yr}. De groeisnelheid is twee en een half keer zo groot
geworden, gelijk op met de uitstoot.
\end{solution}

\begin{solution}{exo:b2:biogeochemical-cycles:7}
\qty{6}{ppm} is \qty{13}{GtC}, tegenover een netto primaire productie van
\qty{35}{GtC}: de zaagtand is een derde daarvan. Hij is kleiner omdat
ademhaling en afbraak de hele zomer doorgaan (de zaagtand registreert de netto
opname, niet de bruto), omdat de seizoenen van het zuidelijk halfrond
tegengesteld zijn en deels compenseren, en omdat de oceaan de schommeling
dempt.
\end{solution}

\begin{solution}{exo:b2:biogeochemical-cycles:8}
$\qty{200}{kg}/(\qty{28}{g/mol}) = 7100$~mol $\mathrm{N_2}$; $16\times
7100 = \num{114000}$~mol ATP; $/30 = 3800$~mol glucose, \qty{690}{kg}.
Een gewas dat \qty{10}{t} droge stof opbouwt, legt ruwweg
\qtyrange{15}{20}{t} suiker vast als de ademhaling wordt meegeteld, zodat
binding het zo'n \qty{4}{\%} van zijn fotosyntheseproducten kost --- de
prijs van onafhankelijkheid van de stikstof in de bodem.
\end{solution}

\begin{solution}{exo:b2:biogeochemical-cycles:9}
Beschikbare verhouding $\mathrm{N:P} = 32/2.2 = 14.5$, onder de 16 van
Redfield: nitraat raakt eerst op, met \qty{0.2}{\micro mol/kg}
fosfaat over. Gebonden koolstof: $32\times 106/16 =
\qty{212}{\micro mol/kg}$, ongeveer \qty{2.5}{mg} koolstof per
kilogram water.
\end{solution}

\begin{solution}{exo:b2:biogeochemical-cycles:10}
$E^{*} = Q\tau = \qty{500}{GtC}$, \qty{235}{ppm} boven het
pre-industriële niveau van 285: ongeveer \qty{520}{ppm}. In werkelijkheid
slechter omdat de putten geen enkele snelle voorraad zijn: de oppervlakte-oceaan
raakt verzadigd (haar vermogen om $\mathrm{CO_2}$ op te nemen daalt naarmate haar
carbonaat wordt verbruikt), de koolstofput op land verzwakt bij opwarming, en
de echte verwijdering heeft een trage component --- menging naar de diepe
oceaan en verwering van gesteente --- met tijdconstanten van duizend tot
honderdduizend jaar, zodat een deel van het overschot op geen enkele $\tau$ van
decennia relaxeert.
\end{solution}

\begin{solution}{exo:b2:biogeochemical-cycles:11}
Wat werd gemeten, is de uitwisseling van atmosferisch $\mathrm{CO_2}$ met de
oceaan en de biosfeer, geen verval: het $\mathrm{{}^{14}C}$ van de kernproeven
werd verdund in de veel grotere voorraden van de oppervlakte-oceaan en van de
planten en bodems, die daarna gewone koolstof teruggaven. De halveringstijd van
11 jaar van het overschot (een e-folding-tijd van ongeveer 16 jaar) is de
tijdschaal van die uitwisseling --- de omzet van de atmosferische koolstof ten
opzichte van de oppervlakte-oceaan en het land, langer dan de bruto
\omterm{def:b2:biogeochemical-cycles:reservoir}{verblijftijd} van vier jaar omdat veel van wat vertrekt meteen terugkomt.
\end{solution}

\begin{solution}{exo:b2:biogeochemical-cycles:12}
Stikstof: natuurlijke binding ongeveer 200~Tg per jaar; menselijke binding
(Haber--Bosch 120, gewassen 60, verbranding 30) ongeveer 210: de \omterm{def:b2:biogeochemical-cycles:reservoir}{flux} is
verdubbeld, en de getroffen \omterm{def:b2:biogeochemical-cycles:reservoir}{reservoirs} --- bodems, rivieren, kustzeeën, het
$\mathrm{N_2O}$ van de lucht --- zijn klein en snel, zodat de verandering
binnen decennia zichtbaar wordt, terwijl het \omterm{def:b2:biogeochemical-cycles:reservoir}{reservoir} $\mathrm{N_2}$ onaangetast
blijft. Koolstof: de menselijke uitstoot van 11~GtC per jaar is maar \qty{5}{\%} van
de bruto natuurlijke uitwisseling van ongeveer 200, zodat de \omterm{def:b2:biogeochemical-cycles:reservoir}{flux} een duwtje
krijgt; maar de natuurlijke \omterm{def:b2:biogeochemical-cycles:reservoir}{fluxen} heffen elkaar op en de menselijke niet, en die
heeft het atmosferische \omterm{def:b2:biogeochemical-cycles:reservoir}{reservoir} met \qty{50}{\%} verhoogd --- het kleine
duwtje heeft een groot cumulatief effect op een \omterm{def:b2:biogeochemical-cycles:reservoir}{reservoir} met een trage
uiteindelijke put. Welke kringloop ``meer verstoord'' is, hangt ervan af of
men \omterm{def:b2:biogeochemical-cycles:reservoir}{fluxen} of opstapeling telt.
\end{solution}

\begin{solution}{pb:b2:biogeochemical-cycles:1}
\textbf{1.} $285\times 2.13 = \qty{607}{GtC}$; $412\times 2.13 =
\qty{878}{GtC}$; toename \qty{271}{GtC}.
\textbf{2.} $271/700 = 0.39$ (de fractie over de recente decennia,
met de uitstoot door landgebruik meegeteld, ligt dichter bij $0.45$).
\textbf{3.} De oceaan, ongeveer een kwart van de uitstoot, en het land
(herstel van bossen en bemesting door $\mathrm{CO_2}$), ongeveer een derde; de
rest van de \qty{430}{GtC} die niet in de lucht zit.
\textbf{4.} $175/38000 = \qty{0.5}{\%}$ --- verwaarloosbaar voor de hele
oceaan; maar de opname is tot nu toe vooral in de oppervlaktelaag terechtgekomen
(\qty{900}{GtC}), waar het om een stijging van een tiende of meer gaat, en de
carbonaatchemie zet een kleine stijging van opgelost $\mathrm{CO_2}$ om in een
grotere stijging van de waterstofionconcentratie: de pH aan het oppervlak is
met $0.1$ gedaald, een stijging van de zuurgraad met \qty{30}{\%}.
\textbf{5.} Met $Q = Q_0 e^{t/T}$ probeert men $E = Ae^{t/T}$: $A/T = Q_0 -
A/\tau$, dus $A = Q_0\tau T/(\tau + T)$, en de \omterm{prop:b2:biogeochemical-cycles:carbon}{atmosferische fractie}
$\dot E/Q = \tau/(\tau + T)$ is constant. Een uitstoot die met
\qty{2}{\%} per jaar groeit ($T = 50$), geeft met $\tau = 40$ $0.44$: de
fractie is constant omdat put en bron samen groeien.
\textbf{6.} Met constante $Q$ gaat $E \to Q\tau$ en $\dot E \to 0$: de
\omterm{prop:b2:biogeochemical-cycles:carbon}{atmosferische fractie} daalt naar nul naarmate de putten bijbenen --- in
het model met één box; in werkelijkheid raken de putten verzadigd en daalt de
fractie trager.
\textbf{7.} \qty{271}{GtC}.
\textbf{8.} $E^{*} = 10\times 100 = \qty{1000}{GtC}$: $285 +
1000/2.13 = \qty{755}{ppm}$.
\textbf{9.} $E(t) = 1000 + (271 - 1000)e^{-t/100}$; bij $t = 80$ is
$E = 1000 - 729\times 0.449 = \qty{672}{GtC}$: \qty{600}{ppm}.
\textbf{10.} $Q = 10(1 - t/50)$: particuliere oplossing $E_p = at + b$
met $a = -20$, $b = 3000$ (uit $a = 10 - b/\tau$ en $0 = -0.2 -
a/\tau$); $E = 3000 - 20t + (271 - 3000)e^{-t/100}$. Bij $t = 50$:
$2000 - 2729\times 0.607 = \qty{344}{GtC}$, \qty{447}{ppm}. Daarna neemt $E$
vrij af: bij $t = 80$ is het $344\,e^{-0.3} = \qty{255}{GtC}$,
\qty{405}{ppm}.
\textbf{11.} \qty{600}{ppm} tegenover \qty{405}{ppm}: het tweede
scenario brengt de lucht tegen 2100 bijna terug naar waar ze nu staat.
\textbf{12.} Met één enkele $\tau$ blijven de putten in hetzelfde tempo werken
op een krimpend overschot; in werkelijkheid hebben de snelle putten (de
menglaag, de groeiende bossen) al opgenomen wat ze kunnen, en gebeurt de
resterende verwijdering door trage menging naar de diepe oceaan en verwering,
met een $\tau$ van eeuwen tot millennia, zodat de daling na netto nul veel
trager verloopt --- de concentratie blijft eeuwenlang ruwweg constant in plaats
van te dalen.
\textbf{13.} Gewas \qty{90}{kg}; gedenitrificeerd \qty{37.5}{kg}; uitgespoeld
\qty{22.5}{kg} stikstof per hectare per jaar.
\textbf{14.} $\qty{22.5}{kg}/\qty{14}{g/mol} = 1600$~mol N; koolstof
$1600\times 106/16 = \num{10600}$~mol, \qty{128}{kg}.
\textbf{15.} \num{10600}~mol $\mathrm{O_2}$, \qty{340}{kg}.
\textbf{16.} $\num{10600}/0.25 = \num{42000}$~\unit{m^3}: de uitspoeling van
één hectare kan elk jaar veertigduizend kubieke meter bodemwater van zijn
zuurstof ontdoen --- en de Mississippi ontwatert honderd miljoen hectare.
\textbf{17.} $37.5\times 0.01 = \qty{0.375}{kg}$ stikstof als
$\mathrm{N_2O}$, dat is $0.375\times 44/28 = \qty{0.59}{kg}$ $\mathrm{N_2O}$;
$\times 270 = \qty{160}{kg}$ $\mathrm{CO_2}$-equivalent.
\textbf{18.} $150\times 4 = \qty{600}{kg}$ $\mathrm{CO_2}$: de productie weegt
vier tegen één zwaarder dan het $\mathrm{N_2O}$ van het veld, en samen zijn ze
drie kwart ton per hectare per jaar.
\textbf{19.} $\tau = 10$ jaar; $M^{*} = 5\times 10 = \qty{50}{t}$;
concentratie $50/(5\times 10^{7}) = \qty{1}{g/m^3} =
\qty{1000}{mg/m^3}$.
\textbf{20.} Dertig keer boven de drempel: sterk eutroof.
\textbf{21.} $1/\tau = 0.1 + 0.2 = 0.3$: $\tau = 3.3$ jaar; $M^{*} =
5/0.3 = \qty{16.7}{t}$, \qty{330}{mg/m^3}.
\textbf{22.} $M^{*} = 1/0.3 = \qty{3.3}{t}$, \qty{67}{mg/m^3} --- nog steeds
eutroof; binnen \qty{10}{\%} na $\tau\ln 10 = 7.7$ jaar.
\textbf{23.} Het sediment is geen put in één richting: fosfor die er in de
eutrofe jaren in is opgeslagen, komt weer vrij onder de anoxische
omstandigheden die de eutrofiëring zelf schept (aan ijzer gebonden
fosfaat lost op wanneer het ijzer wordt gereduceerd), zodat het meer zichzelf
nog jaren voedt nadat de externe aanvoer is gestopt --- interne belasting, een
tweede \omterm{def:b2:biogeochemical-cycles:reservoir}{reservoir} met een lange \omterm{def:b2:biogeochemical-cycles:reservoir}{verblijftijd}.
\textbf{24.} $\qty{5000}{kg}/\qty{31}{g/mol} = 1.6\times 10^{5}$~mol P;
koolstof $\times 106 = 1.7\times 10^{7}$~mol, \qty{205}{t} koolstof per
jaar --- \qty{20}{g/m^2} over het oppervlak van het meer als het
\qty{5}{m} diep is, een zware bloei.
\textbf{25.} \omterm{prop:b2:biogeochemical-cycles:carbon}{Atmosferische fractie} $0.39$ over 1850--2020; in 2100
\qty{600}{ppm} bij een uitstoot die op \qty{10}{GtC/yr} blijft, en
\qty{405}{ppm} bij een uitstoot die tegen 2070 tot nul is teruggebracht (een
optimistisch model met één enkele $\tau$); uitgespoelde stikstof \qty{22.5}{kg}
per hectare per jaar.
\end{solution}
